\section{Reduction of maximal data to a minimal boundary}
\label{sec:maximal-reduction}

The result of this section is a reduction within the asymptotically
hyperbolic class.  In particular, it requires no uniform statement about
the metrics to which the time-symmetric inequality is applied.

\begin{theorem}[Maximal reduction]\label{thm:ah-reduction}
Let \(F:[0,\infty)\to[0,\infty)\) be continuous and nondecreasing.
Suppose that every smooth, complete, connected, orientable, one-ended
asymptotically hyperbolic three-dimensional exterior with nonempty
compact minimal boundary, scalar curvature at least \(-6\),
\(g-b\in C^{2,\gamma}_\tau\) for some \(0<\gamma<1\) and
\(3/2<\tau<3\), the weighted scalar integrability of
Section~\ref{sec:hypotheses}, and four finite mass fluxes satisfies
\[
 p_0\ge F(A_{\min}).
\]
The exponent \(\gamma\) may depend on the comparison metric.
This hypothesis is to hold in every designated asymptotic chart,
without a causal-character assumption on the mass covector.
Then every maximal initial-data exterior of
Section~\ref{sec:hypotheses} satisfies, in each designated chart,
\[
 p_0(g)\ge F\bigl(A_{\min}(S;g)\bigr).
\]
\end{theorem}

Fix the original data and chart, and put \(A_*=A_{\min}(S;g)\).
Extend the end radius \(r\) to a smooth positive function bounded away
from zero on \(\Omega\).  Constants may depend on this extension and
the fixed data.  Derivatives, volume measures, and fluxes in this
section use \(g\), unless a different metric is indicated.
The letter \(b_S\) below denotes a scalar inner flux datum; it is
unrelated to the hyperbolic reference metric \(b\).

\subsection{The equations and the order of choices}

For \(N>0\) and \(f\), define
\begin{equation}\label{eq:ah-fields}
 \begin{gathered}
 s^2+s^4|df|_g^2=N^2,\qquad p=s^2\nabla f,\qquad
 v=p/N,\qquad W=N/s,\\
 \bar g=g+s^2df^2,\qquad
 \pi=N^{-1}\operatorname{sym}\nabla p^\flat,\qquad
 x=sw,\qquad h=x\bar g,\qquad A=Ns\bar g^{-1}.
 \end{gathered}
\end{equation}
The first relation uniquely defines \(s>0\), smoothly also at \(df=0\).
The vector \(p\) in this display is not a component of the mass
covector.  In particular
\(\bar g^{-1}=g^{-1}-v\otimes v\) and
\(\dd V_{\bar g}=W\,\dd V_g\).
We solve
\begin{equation}\label{eq:ah-system}
 \divg_g p=0,\qquad B=\nabla N-Kp,\qquad
 \divg_g B=3N+\Sigma,\qquad
 -\divg_g(A\,dw)=3N+\Sigma x+P ,
\end{equation}
where \(\Sigma,P\ge0\).  At the inner boundary, whose normal
\(\nu\) points into \(\Omega\), prescribe
\begin{equation}\label{eq:ah-inner-data}
 f=tM,\qquad B\cdot\nu=-t b_S,\qquad (A\,dw)\cdot\nu=0 .
\end{equation}
Here \(0\le t\le1\) is a continuation parameter.  On the outer
sphere \(S_R=\{r=R\}\), with outward \(g\)-normal \(\nu_R\), prescribe
\begin{equation}\label{eq:ah-outer-data}
 f=0,\qquad B\cdot\nu_R=\partial_{\nu_R}V_0,\qquad
 w=V_0^{-1}.
\end{equation}
We keep \(K\) fixed throughout the continuation; this preserves
\(H_g(S)\le k\), where \(k=K(\nu,\nu)\).

The identities of Lemmas~\ref{lem:graph-curvature},
\ref{lem:graph-current}, and~\ref{lem:graph-boundary} apply with
\(\tr_gK=\tr_g\pi=0\).  Writing \(Q=B+A\,dw\), they give
\begin{align}
 \divg_g Q&=-(x-1)\Sigma-P,\nonumber\\
 Q&=N\bar\nabla x+x(\pi-K)p-(x-1)B,\label{eq:ah-current}\\
 x^2(\Scal_h+6)
 &=x\{D+|\pi-K|_g^2\}
   +\frac{3}{2x}|dx|_{\bar g}^2+6(x-1)^2+\frac{2P}{N},
 \label{eq:ah-scalar}\\
 D&=16\pi(\mu-J\cdot v)\ge0,\nonumber\\
 \left(\Delta_{\bar g}-3+
  N^{-1}\divg_g((\pi-K)p)\right)x&=-3-P/N .
 \label{eq:ah-x-equation}
\end{align}
The dominant energy condition and \(|v|<1\) give the displayed
sign of \(D\).  At \(S\), the height maximum principle will give
\(v=-z\nu\), \(0\le z<1\), and the boundary formula is
\begin{equation}\label{eq:ah-boundary-curvature}
 s\sqrt{x}\,H_h=N(H_g-zk)-t b_S .
\end{equation}
Consequently a suitable small inner flux will make the obstacle
boundary strictly mean concave wherever an enclosing minimizer
could touch it.

Choose a fixed \(p_1>2\) and an area tolerance \(0<\delta<1/4\).
For a large finite height \(M\), set
\begin{equation}\label{eq:ah-parameter-scales}
 R_s=M^{3/10},\qquad T=M^{4/5},\qquad L=M^2.
\end{equation}
The allowable smooth inner data form the convex class
\begin{equation}\label{eq:ah-flux-class}
 \mathcal B_M=
 \left\{b_S\in C^\infty(S):
 0\le b_S\le B_*,\
 \|b_S\|_{L^{p_1}(S,g)}\le M^{-1/(10p_1)}\right\}.
\end{equation}
The fixed constant \(B_*\) will be chosen using an upper lapse
estimate which depends only on the norm constraint in
\eqref{eq:ah-flux-class}, not on \(B_*\).

Let \(\epsilon>0\) be chosen after \(M,L\).  The unscaled volume
source is a product of smooth cutoffs, between \(0\) and \(L\),
which vanishes when \(r\ge2R_s\), \(x\ge1-\delta/2\), or
\(N\le\epsilon\), and equals \(L\) when
\[
 r\le R_s,\qquad x\le1-\delta,\qquad N\ge2\epsilon .
\]
The smooth cutoffs are defined also for negative arguments.
Two further sources, each bounded by a finite strength \(d_*\),
have the following support and full-strength regions:
\begin{equation}\label{eq:ah-penalty-cutoffs}
 \begin{array}{c|c|c}
 &\text{vanishes if}&\text{equals }d_*\text{ if}\\ \hline
 P_1& f\le M/2\ \text{or }w\ge4T&
       f\ge3M/4,\quad w\le2T\\
 P_2& N\ge4\epsilon\ \text{or }w\ge4T_2&
       N\le3\epsilon,\quad w\le2T_2 .
 \end{array}
\end{equation}
Here \(T_2=C_2(M,L)/\epsilon\), with \(C_2\) fixed below before
\(\epsilon\) is chosen.  During the continuation, multiply both
the volume source and \(P_1+P_2\) by \(t\); we continue to denote
the scaled sources by \(\Sigma\), \(P_1\), \(P_2\),
and \(P=P_1+P_2\).
All strengths and heights remain finite when solving the equations.
The limit \(R\to\infty\) is always taken with these internal
parameters fixed.

\subsection{Lapse estimates with only bounded drift}

\begin{lemma}\label{lem:ah-lapse}
For smooth solutions on sufficiently large truncations, uniformly
over \(0\le t\le1\) and \(b_S\in\mathcal B_M\),
\begin{equation}\label{eq:ah-lapse-bounds}
 N\le Cr,\qquad N\ge c(M)\epsilon r,\qquad
 \frac{|N-V_0|}{r}\le
 C(R_s/r)^{1+\kappa}\quad(r\ge2R_s),
\end{equation}
where \(0<\kappa<\tau-1\) is fixed and small.  The upper constant
uses only the \(L^{p_1}\) bound of \(b_S\).  On compact sets,
\(N\) has uniform Hölder and \(W^{1,q}\) bounds for some \(q>3\),
depending on \(M,L\), but independent of \(\epsilon,d_*\) and
derivatives of \(b_S\).

The positivity and lower bound remain true at a fixed point in
which the drift \(Z=K(v)\) is computed using a smooth positive
extension of the constitutive lapse argument, provided that the
actual lapse equation is
\(\divg_g(\nabla N-ZN)=3N+\Sigma\) and the cutoff
\(N\le\epsilon\) uses the actual lapse.
\end{lemma}

\begin{proof}
Only \(|Z|\le |K|\) is used.  First consider the density equation
\begin{equation}\label{eq:ah-homogeneous-lapse}
 \divg_g(\nabla H_1-ZH_1)=3H_1
\end{equation}
with the prescribed outer flux and inner outward flux \(t b_S\).
Its adjoint is \(\Delta_g+Z\cdot\nabla-3\), with ordinary
homogeneous Neumann boundary condition.  The maximum and boundary
point principles give zero adjoint kernel.  Scalar elliptic
Fredholm theory, deforming the lower terms to those of
\(\Delta_g-3\), gives index zero and hence unique solvability of
the primal problem.  Green's identity with the adjoint solution
of a nonnegative forcing, which is nonpositive, proves
\(H_1\ge0\): both prescribed outward fluxes are nonnegative, and
the outer one is strictly positive.  Local Harnack then gives
strict positivity.

We give the weighted estimate used here, including its uniformity
in the outer radius.  For \(-\Delta_g+3\), scalar and divergence
forcing of size \(O(r^{-j})\), \(0<j<3\), with natural flux data
in \(L^{p_1}(S)\) and of size \(O(R^{-j})\) on \(S_R\), yield
an \(O(r^{-j})\) solution, with a constant independent of \(R\).
The natural datum includes the flux of the divergence forcing.
To see this, construct a positive strict supersolution comparable
to \(r^{-j}\), constant on a fixed inner region.  Change its
logarithmic slope slowly from \(0\) to \(-j\) in a long interval
of \(\log r\) lying sufficiently far out.  The limiting radial
indicial slopes are \(-3\) and \(1\), so this retains
\((-\Delta_g+3)\varphi\ge c\varphi\).  Add a fixed sufficiently
large multiple of
\(R^{-j-a}r^a\), similarly continued as a constant inside, where
\(0<a<1\), to make its outer normal derivative nonnegative.
The resulting barrier stays comparable to \(r^{-j}\) on \(r\le R\).

For divergence forcing the barrier argument is justified by
normalization and compactness.  If the asserted bound fails,
divide a violating solution by the maximum of its absolute
value relative to the barrier, and normalize locally by the
barrier's size at a maximizing point.  The normalized forcing
and boundary data tend to zero.  Local Laplace \(W^{1,q}\)
estimates are available with
\(3<q<3p_1/2\): the boundary \(L^{p_1}\) datum belongs to the
required negative trace space by the two-dimensional boundary
Sobolev embedding.  These estimates give local uniform compactness.
The outer sphere charts have uniformly controlled geometry.
They may be formed by flowing
\(\nabla\eta/|d\eta|^2\), \(\eta=\operatorname{arsinh}r\);
the normal direction is orthogonal and the reflected principal
metric is uniformly Lipschitz.  Thus no high derivatives of
geodesic normal coordinates are needed.
At points escaping to infinity, subsequences of the normalized
geometries converge to hyperbolic geometry, the relative radius
to an exponential horospherical coordinate, and a nearby outer
boundary to a horosphere.  The strict supersolution and its
Neumann sign persist in these limits.  A nonzero maximum of the
limiting solution divided by that supersolution contradicts the
maximum or boundary point principle.  The same argument in a
fixed chart handles bounded maximizing points.  This proves the
claimed estimate.

Apply it to \(H_1-V_0\), retaining the drift as divergence forcing.
The background errors are \(O(r^{1-\tau})\), and
\(Z=O(r^{-\tau})\).  It follows that
\begin{equation}\label{eq:ah-homogeneous-asymptotic}
 |H_1-V_0|\le Cr^{-\kappa}.
\end{equation}
Here and below \(V_0\) is extended smoothly across the compact
part.  There is one additional compactness point in this
application: a failure with bounded maximizing points could
produce a nonzero limit \(H_\circ\ge0\), with bounded measurable
limit drift, zero natural flux, and
\(H_\circ=O(r^{-\kappa})\), solving
\eqref{eq:ah-homogeneous-lapse}.  The sign follows from positivity
of \(H_1\) before subtracting \(V_0\) and dividing by an unbounded
normalizing constant.  Regard \(ZH_\circ\) as divergence forcing
for \(-\Delta_g+3\) and apply the preceding estimate repeatedly.
This improves the decay to \(O(r^{-j'})\) for some \(2<j'<3\).
To justify each improvement on the full exterior, solve the
baseline problem on truncations with zero natural outer flux
and pass to a limit; its difference from \(H_\circ\) decays in
the previous weight and vanishes by the maximum principle.
The improved decay is integrable.  Testing the density equation
with radial cutoffs, equal to one near the finite boundary,
and moving the Laplacian onto the cutoffs gives
\(3\int_\Omega H_\circ=0\).  All shell terms vanish since
\(j'>2\).  This is impossible.  On bounded limiting truncations
the same contradiction follows by integrating the zero-flux
equation directly.  This proves
\eqref{eq:ah-homogeneous-asymptotic} with the asserted uniformity.
For \(b_S=0\), Harnack, also by homogeneous-flux reflection at
\(S\), and connectedness now give \(H_1\ge cr\).

The dual maximum principle compares \(N\) to the homogeneous
solution with the same drift and boundary data, giving
\(N\le H_1\le Cr\).  For the lower bound divide \(N\) by the
homogeneous solution with the same drift and \(b_S=0\).
If this ratio is below \(c\epsilon/R_s\), with \(c\) sufficiently
small, then \(N\le\epsilon\) wherever the spatial cutoff permits
\(\Sigma\ne0\).  Hence \(\Sigma=0\) on this sublevel.  The ratio
there solves an elliptic drift equation with no zeroth term.
Its outward derivative is nonnegative at the inner face and
strictly positive at the outer face when the ratio is below
the threshold.  The minimum and boundary point principles
therefore exclude such a minimum.  This proves the positive
floor, and also excludes negative lapses when constitutive
arguments have been extended.

Finally, on \(2R_s\le r\le R\), subtract \(V_0\) and treat \(ZN\)
as divergence forcing of size \(O(r^{1-\tau})\).
The inner Dirichlet values have size \(O(R_s)\), and the natural
outer error has the corresponding decaying size.  The weighted
argument above, now with the inner Dirichlet face, gives
\[
 |N-V_0|\le C R_s^{1+\kappa}r^{-\kappa}.
\]
The homogeneous contribution of the inner data is bounded by
the same barriers.  In the normalized argument the zero
Dirichlet trace also excludes a limit maximum at that face.
This proves the final estimate in \eqref{eq:ah-lapse-bounds}.
On fixed patches write the lapse equation as a metric Laplace
equation with bounded divergence forcing \(ZN\), bounded scalar
forcing \(3N+\Sigma\), and the prescribed \(L^{p_1}\) boundary
data.  The local \(W^{1,q}\) and Hölder estimates from
Proposition~\ref{prop:elliptic-inputs} prove the last assertion.
\end{proof}

\subsection{Finite slopes and uniform estimates at infinity}

\begin{lemma}\label{lem:ah-graph-estimates}
The solutions satisfy \(0\le f\le tM\).  On every fixed compact
subset, including \(S\),
\begin{equation}\label{eq:ah-s-floor}
 s\ge c_{\mathrm{loc}}(M,L)\epsilon ,
\end{equation}
and their height gradients have finite bounds for fixed
\(M,L,\epsilon\), independent of \(d_*\), \(R\), and derivatives
of \(b_S\).  Moreover,
\begin{align}
 \int_{\{f\in I\}}p\cdot df&\le C|I|,\qquad
 \int_{\Omega_R}(W^2-1)\,\dd V_g\le CM,\label{eq:ah-height-energy}\\
 |f|&\le C R_s^{-1}(R_s/r)^{4-\xi},\nonumber\\
 |v|+|\pi|+|\nabla\pi|&\le C(R_s/r)^{3-\xi}
       \quad(r\ge C_0R_s),\label{eq:ah-end-graph}
\end{align}
where \(I\) is any height interval, \(\xi>0\) can be fixed
sufficiently small, and the constants in
\eqref{eq:ah-height-energy}--\eqref{eq:ah-end-graph} are
independent of \(M,L,\epsilon,d_*\) for large \(M\).
The end estimates include the corresponding local Hölder
control and derivatives of \(v\) through order two.
\end{lemma}

\begin{proof}
The maximum principle for \(\divg_g p=0\) gives the height bound.
We spell out how the differential estimate in
Lemma~\ref{lem:gradient-maximum} supplies all the finite bounds
used here.  At nonzero gradient write
\[
 z=|v|,\quad y=z^2,\quad m_f=|df|,\quad
 q=\frac{1-y}{1+y},\quad
 \beta=1+q,\quad
 \mathcal P=g^{-1}+(q-1)e\otimes e,\quad e=\nabla f/m_f .
\]
The height equation is
\(\mathcal P:\nabla^2 f+\beta\,d\log N\cdot df=0\).
At an interior maximum of
\(G+j(f)+d_0\), where \(G=-\log s\), \(j'<0\),
\(j''>0\), and \(y\ge1/2\), that lemma gives
\begin{equation}\label{eq:ah-gradient-maximum}
 j''q m_f^2\le
 C\left[1+\frac{\Sigma}{N}
 +(1-y)^2(j'm_f)^2+|dd_0|^2+|\nabla^2d_0|\right].
\end{equation}
Its hypotheses hold because the lapse equation has precisely
the drift \(K(v)\), the tensor \(K\) and its divergence are
bounded on the relevant patches, and \(\Sigma\ge0\).
The cancellation of the Hessian of \(\log N\) in that lemma
requires no prior second-derivative estimate for \(N\).
Multiplying \eqref{eq:ah-gradient-maximum} by \(N^2\) controls
\((1-y)^{-1}\), using
\(Nm_f=z/(1-y)\) and the upper lapse bound; a positive lower
lapse is not needed in this step.

At either constant-height face let \(n\) point towards increasing
end direction.  When \(df\ne0\), \(e=-n\), and the boundary part
of the same lemma gives
\[
 G_n=-\frac{yH_g+\partial_n\log N}{1-y}.
\]
At \(S\), \(\partial_\nu\log N=-zk-tb_S/N\), so
\(G_\nu\ge zk/(1+z)\), a uniform lower bound.
Since \(j'f_\nu=-j'm_f>0\), a maximum at this face bounds \(m_f\).
At \(S_R\), \(H_g\) is bounded below by a positive constant,
\(K=o(1)\), and
\(N\partial_{\nu_R}\log N=K(p,\nu_R)+\partial_{\nu_R}V_0\).
Taking \(|j'|\) sufficiently small makes the outward derivative
of \(G+j(f)\) negative at large \(y\), bounded above by
\(-c/(1-y)\).  Thus an outer maximum is excluded, or its slope
is bounded by a localization derivative.  A sufficiently small
decreasing exponential \(j\) has both required derivative signs.

On a compact truncation these arguments give a finite global
slope bound.  To make it independent of \(R\) locally, set
\(d_0=k_0\log\chi\), \(k_0\ge2\), where \(\chi\) is a smooth
cutoff on a slightly larger compact region, constant normally
near \(S\) if that face is present.  Equation
\eqref{eq:ah-gradient-maximum} bounds \(W\chi\) at an interior
maximum; the preceding face argument does the same at \(S\).
When \(y<1/2\) this bound is immediate.  Applying the lapse floor
to the weighted maximum gives \eqref{eq:ah-s-floor}.

Since \(|p|\le N\), the flux through \(S\) is bounded independently
of all internal large parameters.  Test \(\divg_g p=0\) with
a clipped height whose derivative is one on \(I\) and zero
elsewhere.  Its boundary oscillation is at most \(|I|\), giving
the first estimate in \eqref{eq:ah-height-energy}.
The identity \(p\cdot df=W^2-1\), tested with the entire height,
gives the second one.

We next prove the initial uniform end bound \(f\le C/r\);
it cannot be obtained from a compact-region bound depending
arbitrarily on \(M\).  By Lemma~\ref{lem:ah-lapse}, \(N\asymp r\)
when \(r\ge CR_s\), for a fixed large \(C\).
Fix a radial segment of \(\eta\)-length at most one starting
at radius \(r_*\) there, and set \(F=r_*f\).  In a bounded
number of uniformly controlled metric charts covering its
neighborhood, consider the subgraph of \(F\) in the product
with a Euclidean line.  The product vector
\[
 \mathcal Z=(-p/r_*,1)
\]
is bounded and divergence free.  Its pairing with the upward
graph normal is
\[
 \frac{1+p\cdot df}{\sqrt{1+r_*^2|df|^2}}
 =\frac{W^2}{\sqrt{1+r_*^2|df|^2}}\ge c>0,
\]
using \(N/r_*\asymp1\) and the constitutive relation.
The divergence theorem therefore makes the subgraph locally
quasiminimizing for perimeter, with a uniform constant.
Filling or emptying a product ball bounds the graph perimeter
there by a constant times the trace area of either phase.
The four-dimensional isoperimetric inequality then gives,
for the positive volume \(V(h)\) of either phase in a ball,
\(V'(h)\ge cV(h)^{3/4}\).
Integration yields uniform two-sided volume density, and
relative isoperimetry yields a fixed positive graph-area
lower bound in each sufficiently small ball centered on
the graph.

If the height span of \(F\) on the radial segment is \(D\),
choose graph points with separated heights and disjoint
fixed-size product balls.  The preceding estimate bounds
the graph area in the surrounding tube and relevant height
band below by \(cD-C\).  Conversely, the constitutive relation
gives
\(\sqrt{1+r_*^2|df|^2}\le C(1+p\cdot df)\).
The base tube has bounded volume and the relevant interval
of \(f\) has length at most \(C(D+1)/r_*\).
Equation~\eqref{eq:ah-height-energy} bounds the graph area
above by \(C+C(D+1)/r_*\).  For large fixed starting radius
this bounds \(D\) uniformly.  If the segment meets \(S_R\),
oddly reflect \(f\), evenly reflect \(N\), and reflect the
metric in the orthogonal charts described in the lapse proof.
The zero trace of \(f\) makes the divergence equation valid
weakly after reflection.  The reflected height band adds
only an interval of comparable length to the energy estimate.
Thus the same estimate holds there.  Summing the bounds
\(C/r_*\) over consecutive outward unit segments, up to
the zero height at \(S_R\), proves \(f\le C/r\).

On far-out unit balls rescale \((f,N,p,s)\) by
\((r_*,r_*^{-1},r_*^{-1},r_*^{-1})\), respectively.
The equations there have \(\Sigma=0\) and retain their form.
The localized gradient estimate now gives uniform bounds
for the rescaled height gradient and uniform ellipticity.
The outer-face argument applies with its rescaled positive
flux datum.  The rescaled lapse has uniform \(W^{1,q}\) bounds.
Differentiating the divergence height equation, or using
difference quotients, gives scalar uniformly elliptic
divergence equations for its first derivatives, with divergence
forcing in \(L^q\), \(q>3\).  The inhomogeneous Hölder estimate
therefore bounds \(df\) in a Hölder norm, also at a
constant-height boundary by odd reflection.  The lapse
divergence equation then gives \(C^{1,\gamma}\) control,
followed by the nondivergence height equation and successive
Schauder estimates.  These steps use only the stated
\(C^{2,\alpha}\) metric and \(C^{1,\alpha}\) tensor bounds.
They also prove compact-set gradient Hölder estimates away
from infinity without differentiating \(b_S\).

On each far-out ball, the energy estimate for a height interval
of length \(C/r_*\), together with this Hölder control, implies
\(z=o(1)\) uniformly as \(r_*\to\infty\).
The derivative version of \eqref{eq:ah-lapse-bounds} implies
that \(\nabla\log N\) tends to the hyperbolic radial direction
as \(r/R_s\to\infty\).  The limiting height operator is
\(\Delta_b+2\,d\eta\cdot d\); its decaying radial indicial
root is \(-4\).  For any sufficiently small fixed \(\xi>0\),
a positive multiple of \(r^{-4+\xi}\) is therefore a strict
supersolution sufficiently far beyond \(R_s\).
Comparison with the already proved \(C/r\) bound at a fixed
multiple of \(R_s\) gives the first estimate in
\eqref{eq:ah-end-graph}.  Treating the actual height equation
as a linear equation on these smaller height ranges gives
the corresponding derivative bounds by Schauder theory.
Inserting them in \(p=s^2\nabla f\) and
\(\pi=N^{-1}\operatorname{sym}\nabla p^\flat\) gives the
remaining assertions.
\end{proof}

For a fixed sufficiently large \(C_0\), all sources vanish when
\(r\ge C_0R_s\).  Indeed the height is then below \(M/2\),
the lapse exceeds \(4\epsilon\), and the spatial cutoff
for \(\Sigma\) has ended.  In this region
\begin{equation}\label{eq:ah-end-x}
 (\Delta_{\bar g}-3-U_0)x=-3,\qquad
 U_0=-N^{-1}\divg_g((\pi-K)p),\qquad
 |U_0|_{C^{0,\gamma}_{\mathrm{loc}}}
       \le C(R_s/r)^\lambda
\end{equation}
for some \(\lambda>3\).  For example one can take any exponent
below \(\min(6-2\xi,\tau+3-\xi)\), by
\eqref{eq:ah-end-graph}.  The constants here retain the
independence asserted in that lemma.

\subsection{Costs and enforcement of the finite penalties}

\begin{lemma}\label{lem:ah-penalties}
The parameters can be chosen in the order already specified,
with \(\epsilon\) sufficiently small after \(M,L\), and then
\(d_*\) sufficiently large but finite, such that
\begin{equation}\label{eq:ah-penalty-costs}
 \int_{\Omega_R}P_1\,\dd V_g\le C T/M,\qquad
 \int_{\Omega_R}P_2\,\dd V_g=o_{M,L}(1)
       \quad(\epsilon\downarrow0),
\end{equation}
uniformly in \(R,b_S\), and in the strengths.
At \(t=1\), the same choices enforce
\begin{equation}\label{eq:ah-penalty-targets}
 w|_S\ge T,\qquad
 w\ge T_2\quad\hbox{on }\{N\le2\epsilon\},\qquad
 sT_2\ge1\quad\hbox{on }\{N\le4\epsilon\}.
\end{equation}
Consequently \(x<1-\delta\) and \(r\le R_s\) imply
\(\Sigma=L\).
\end{lemma}

\begin{proof}
The lapse floor confines \(N<6\epsilon\) to a compact set
depending on \(M\), independently of \(\epsilon\).
Equation~\eqref{eq:ah-s-floor} on a slightly larger compact set
therefore permits a choice of \(C_2(M,L)\) such that
\(sT_2\ge1\) throughout the required low-lapse region.

For every Lipschitz test \(\zeta\) vanishing at \(S_R\), positivity
of \(w\), its equation, and its zero inner flux give
\begin{equation}\label{eq:ah-penalty-test}
 \int P\,\frac{\zeta^2}{w}\,\dd V_g
       \le \int A(d\zeta,d\zeta)\,\dd V_g .
\end{equation}
Indeed test with \(\zeta^2/w\) and complete the square in
\(2\zeta\,dw\cdot A\,d\zeta/w-\zeta^2 A(dw,dw)/w^2\).
All the other terms on the right side of the \(w\)-equation
are nonnegative.  On the support of \(P_1\),
\(f/M\ge1/2\) and \(w\le4T\).  With \(\zeta=f/M\),
the algebraic inequality
\(A(df,df)=W-W^{-1}\le W^2-1=p\cdot df\) and
\eqref{eq:ah-height-energy} prove
\(\int P_1\le CT/M\).

For \(P_2\) take a cutoff of \(N/\epsilon\), equal to one
below \(4\) and zero above \(6\).  It vanishes near \(S_R\)
for large \(R\).  Since \(A\le N^2g^{-1}\) and \(w\le4T_2\)
on this source's support,
\begin{equation}\label{eq:ah-low-lapse-cost}
 \int P_2\le
 C T_2\int_{\{N<6\epsilon\}}|\nabla N|^2\,\dd V_g.
\end{equation}
Test the lapse equation with \((6\epsilon-N)^+\).
The outward current flux at \(S\) is \(t b_S\ge0\), so its
boundary contribution has the favorable sign.  On the
sublevel set \(|Kp|\le CN\le C\epsilon\).  Cauchy's inequality
then gives
\begin{equation}\label{eq:ah-low-lapse-energy}
 \int_{\{N<6\epsilon\}}|\nabla N|^2
 \le C(M)\epsilon^2
       +C\epsilon\int_{\{N<6\epsilon\}}\Sigma .
\end{equation}

We must show that the final integral tends uniformly to zero;
boundedness of \(\Sigma\) alone would not suffice.
For any potentially violating sequence, use the uniform
local lapse Hölder bounds to take \(N_j\to N_0\ge0\)
uniformly on the compact region containing these sublevels.
After a subsequence,
\[
 Kp_j\rightharpoonup Z_0N_0,\qquad
 \Sigma_j\rightharpoonup^*\Sigma_0,\qquad
 |Z_0|\le C,\quad 0\le\Sigma_0\le L .
\]
The weak limit equation is
\(\Delta_gN_0=\divg_g(Z_0N_0)+3N_0+\Sigma_0\).
We claim \(\Sigma_0=0\) almost everywhere on \(\{N_0=0\}\).
Choose an interior density point of that zero set which is
also a Lebesgue point of \(\Sigma_0\).  In balls of radius
\(h\downarrow0\), rescale \(N_0\) by \(h^{-2}\).
The rescaled fields are nonnegative and vanish at the center.
Their divergence drift has size \(O(h)\) and their scalar
forcing is bounded.  The inhomogeneous Harnack estimate
therefore gives a uniform bound on smaller balls.
Because the rescaled zero sets have asymptotically full
measure, these functions tend to zero in local \(L^1\).
Testing the rescaled equation against compactly supported
smooth functions shows that the limiting constant scalar
forcing is zero, proving the claim.  Boundary points need
not be considered, since \(S\) has zero volume.

For completeness, the changing sublevel sets cause no loss:
for every \(a>0\), uniform convergence gives
\(\{N_j<6\epsilon_j\}\subset\{N_0\le a\}\) eventually.
Nonnegativity and weak convergence, using continuous cutoffs
of \(N_0\) that majorize this fixed sublevel, imply
\[
 \limsup_j\int_{\{N_j<6\epsilon_j\}}\Sigma_j
 \le \int_{\{N_0\le2a\}}\Sigma_0.
\]
Letting \(a\downarrow0\) proves the asserted uniform convergence.
Equations~\eqref{eq:ah-low-lapse-cost} and
\eqref{eq:ah-low-lapse-energy}, with \(T_2=C_2/\epsilon\),
prove the second cost in \eqref{eq:ah-penalty-costs}.
Choose \(\epsilon=\epsilon(M)\) sufficiently small that this
cost, as well as \(T/M=M^{-1/5}\), tends to zero as \(M\to\infty\).

Now fix \(M,L,\epsilon\) and increase the finite strengths.
The previously proved continuity of \(f,N\) supplies uniform
neighborhoods of the targets \(S\) and \(N\le2\epsilon\)
inside, respectively, the weaker triggering regions
\(f\ge3M/4\) and \(N\le3\epsilon\).
The cost bounds imply that the volume in these neighborhoods
where \(w\le2T\), or \(w\le2T_2\), tends to zero as
\(d_*\to\infty\).
For either threshold \(D=T,T_2\), the bounded nonnegative
function \((2D-w)^+\) is a weak subsolution of the homogeneous
divergence operator \(\divg_g(A\,d\cdot)\).
Its coefficients are uniformly elliptic on the relevant
compact region for these fixed parameters, independently of
the strength.  The scalar local supremum estimate bounds its
supremum on a smaller neighborhood by its \(L^2\) mean on the
larger one.  At \(S\), even reflection using the zero natural
flux gives the same estimate.  Choosing \(d_*\) sufficiently
large makes that supremum smaller than \(D\).
This proves \eqref{eq:ah-penalty-targets}, with room in the
thresholds if desired.  Finally, if \(x<1-\delta\), the
second target excludes \(N\le2\epsilon\), so the volume
penalty is at full strength when also \(r\le R_s\).
\end{proof}

\subsection{Compactness, mass, and a uniform comparison sphere}

Until the existence argument below, a \emph{limit solution} means
any limit along expanding truncations of smooth solutions at
the fixed internal parameters and fixed \(t\), allowing arbitrary
\(b_S\in\mathcal B_M\).  Statements about this family are a priori
statements and do not presuppose that it is nonempty.
Pass also to a weak \(L^{p_1}(S)\) limit of the boundary data.
Until the selector is imposed, we use \(b_S\) for this inherited
datum in a limit solution: it need not be smooth, but its
pointwise bounds, norm bound, and natural-flux identity pass
to the limit.  Its integral is the limit of the integrals of
the smooth data.

\begin{lemma}\label{lem:ah-limit-mass}
With internal parameters fixed, the fields \(N,f,df,w\)
are uniformly precompact in local uniform topology, including
at \(S\).  Limit solutions have \(N,w>0\) and are classical
away from \(S\).  They obey
\begin{equation}\label{eq:ah-conformal-decay}
 x-1=O_{\mathrm{fixed},2,\gamma}(r^{-3+\xi}).
\end{equation}
Here \(\xi\) is taken sufficiently small, in particular
\(0<\xi<\min(3-\tau,\tau-1)\).
Their metrics \(h\) satisfy the required asymptotically
hyperbolic decay, all four mass fluxes are finite, and
\(V_0(\Scal_h+6)\) is integrable near infinity.  Their time
components satisfy
\begin{equation}\label{eq:ah-mass-budget}
 8\pi\{p_0(h)-p_0(g)\}
   =t\int_S b_S\,\dd A_g
      +\int_\Omega\{(x-1)\Sigma+P\}\,\dd V_g
   \le o(1)\quad(M\to\infty).
\end{equation}
The \(o(1)\) is uniform over the allowed data, with the
subsequent choices of \(\epsilon,d_*\) just described.
\end{lemma}

\begin{proof}
Only the upper bound for \(w\) remains to be established.
The right side of its equation is bounded on compact sets:
\(\Sigma x\le L\), since a nonzero \(\Sigma\) requires \(x<1\),
and \(P\) has fixed finite strength.
By \eqref{eq:ah-lapse-bounds} and \eqref{eq:ah-end-graph},
a multiple of \(r^{-1}\) is a supersolution for
\(-\divg_g(A\,dw)=3N\) sufficiently far beyond \(R_s\).
Comparison therefore bounds the end values of \(w\) by
\[
 C\left(1+\sup_{\{r\le C_0R_s\}}w\right)r^{-1}.
\]
If its compact supremum were unbounded, divide by that
supremum.  Uniform ellipticity, local divergence estimates,
and convergence of \(N,df\) then give a nonzero nonnegative
limit satisfying the homogeneous \(w\)-equation, with zero
inner flux and decaying to zero at infinity.  The maximum
principle excludes such a function; Harnack and reflection
at \(S\) handle a maximum at that face.  Thus the compact
upper bounds hold.  Local Hölder estimates for \(w\) and
the already proved estimates for \(N,df\) give the stated
precompactness.  The lapse floor and the strictly positive
forcing \(3N\) in the supersolution equation imply positivity
of both limiting fields.  Scalar elliptic bootstrapping gives
classical interior regularity.

On the end, \eqref{eq:ah-end-x} and these upper bounds permit
comparison for \(x-1\) with a multiple of \(r^{-3+\xi}\).
For a finite truncation add \(C(r/R)^{a_0}\), \(0<a_0<1\),
to dominate its outer boundary values.  The first term
dominates the decaying forcing \(U_0\), while the second
vanishes on fixed sets as \(R\to\infty\).
The scalar zeroth coefficient is positive after increasing
the fixed end radius.  The maximum principle proves the
zeroth-order bound in \eqref{eq:ah-conformal-decay}, and
interior Schauder estimates prove its derivative and Hölder
versions.  These estimates are uniform over all limit
solutions at the fixed parameters, without bounds on
derivatives of \(b_S\).

Since
\(\bar g-g=O_{\mathrm{fixed},2,\gamma}(r^{-6+2\xi})\),
the choice \(\xi<3-\tau\) ensures that \(h-b\) has the
original admissible decay exponent, with a possibly smaller
Hölder exponent.  Put \(\varphi=x-1\).
The change in each defining mass flux is, up to an error
tending to zero, the flux of the one-form
\[
 -2\{V_\mu\,d\varphi-\varphi\,dV_\mu\}.
\]
All nonlinear conformal errors and graph terms disappear
by the displayed decay.  This flux has a finite limit:
its divergence is
\(-2V_\mu(\Delta_b-3)\varphi\), and
\eqref{eq:ah-end-x} gives absolute weighted integrability.
Indeed \(U_0=O(r^{-\lambda})\), \(\lambda>3\);
the difference between \(\Delta_{\bar g}\) and \(\Delta_b\)
on \(\varphi\) is
\(O(r^{-\tau-3+\xi})+O(r^{-9+3\xi})\).
Multiplication by \(V_0\) and integration against the
hyperbolic volume, of radial size \(r\,\dd r\), is integrable.
The scalar identity \eqref{eq:ah-scalar}, the original
weighted matter and \(K^2\) integrability, and
\eqref{eq:ah-end-graph}--\eqref{eq:ah-conformal-decay}
similarly prove integrability of \(V_0(\Scal_h+6)\).

On each limit end the lapse equation also gives
\[
 (\Delta_g-3)(N-V_0)
  =O(r^{1-\tau})+\divg_g(Kp).
\]
Together with \eqref{eq:ah-lapse-bounds} and the end
derivative estimates this yields
\(|\nabla(N-V_0)|=O_{\mathrm{fixed}}(r^{-\kappa})\).
Insert the second expression for \(Q\) in
\eqref{eq:ah-current}.  Its sphere flux differs from
that of \(V_0\,d\varphi-\varphi\,dV_0\) by a quantity
tending to zero.  The \(N-V_0\) terms have integrated
size \(O(r^{-1-\kappa+\xi})\); the
\((\pi-K)p\) terms have integrated sizes bounded by
\(O(r^{-3+2\xi})+O(r^{-\tau+\xi})\).
Metric replacement errors vanish by the same decay.
It follows that
\[
 \lim_{r\to\infty}\int_{\{r=\mathrm{const}\}}Q\cdot\nu_r
       \,\dd A_g=-8\pi\{p_0(h)-p_0(g)\}.
\]
At the inner boundary the outward normal of the exterior
is \(-\nu\), and \(Q\cdot(-\nu)=t b_S\).
Integrating \(\divg_g Q=-(x-1)\Sigma-P\) proves
\eqref{eq:ah-mass-budget}.  This integration remains valid
for limits with only continuous fields at \(S\): use the
first expression \(Q=A\,dw+B\), the weak equations and
their prescribed natural fluxes on compact regions.
The second expression is needed only at infinity, where
the fields are classical.  Finally,
\((x-1)\Sigma\le0\), the \(L^{p_1}\) bound makes
\(\int_S b_S=o(1)\), and
Lemma~\ref{lem:ah-penalties} gives \(\int P=o(1)\).
\end{proof}

The next estimate is needed before choosing \(b_S\); compact
upper bounds for \(w\) at fixed parameters alone would not
give it.

\begin{lemma}\label{lem:ah-comparison-sphere}
There is a fixed large \(C_1\) such that, at
\(r_1=C_1R_s\), every limit solution at \(t=1\) satisfies
\begin{equation}\label{eq:ah-comparison-area}
 \Area_h(\{r=r_1\})\le C R_s^2,
\end{equation}
with \(C\) independent of the internal large parameters.
\end{lemma}

\begin{proof}
On \(r>r_1\) solve
\[
 (\Delta_{\bar g}-3-U_0)Y=0,\qquad
 Y|_{r=r_1}=0,\qquad Y/V_0\longrightarrow1 .
\]
Existence follows by solving with outer Dirichlet datum
\(V_0\) on larger truncations, applying positive barriers,
and passing to a limit.  Choose
\(\xi<b_0<\min(1,\tau)\).
Positive and negative barriers of size
\(Cr(r_1/r)^{b_0}\) for \(Y-V_0\) work by the end estimates.
They give \(0\le Y\le Cr\), and the corresponding first
derivative control follows from elliptic estimates.
On an annulus from \(r_1\) to a large fixed multiple of
\(r_1\), compare from below with a small positive multiple
of \(r_1((r/r_1)^2-1)\).  This is a subsolution with zero
inner value, and its outer value can be made smaller than
that of \(Y\).  The boundary point principle and local
boundary estimates therefore give, in the end direction,
\[
 cr_1\le\partial_nY\le Cr_1
      \quad\hbox{on }\{r=r_1\},
\]
where \(n\) is the \(\bar g\)-unit normal.

For \(\varphi=x-1\), the operator just used satisfies
\((\Delta_{\bar g}-3-U_0)\varphi=U_0\).
Green's identity, the mass-flux calculation, and
\(b_0>\xi\) give
\begin{equation}\label{eq:ah-green-comparison}
 \int_{\{r=r_1\}}(x-1)\partial_nY\,\dd A_{\bar g}
 =\int_{\{r>r_1\}}YU_0\,\dd V_{\bar g}
       +8\pi\{p_0(h)-p_0(g)\}
 \le Cr_1^3 .
\end{equation}
The source integral is bounded by \(Cr_1^3\), since
\(\lambda>3\), \(Y\le Cr\), and \(W\) is uniformly bounded
on this end region.  The mass increment is bounded above
by \eqref{eq:ah-mass-budget}.  Add
\(\int\partial_nY\le Cr_1^3\) to both sides of
\eqref{eq:ah-green-comparison}, and then use \(x>0\) and
\(\partial_nY\ge cr_1\).  This yields
\(\int_{\{r=r_1\}}x\,\dd A_{\bar g}\le Cr_1^2\).
Two-dimensional conformal scaling makes this integral
exactly the area in \eqref{eq:ah-comparison-area}.
\end{proof}

\subsection{A continuous choice of inner flux}

\begin{lemma}\label{lem:ah-contact-selector}
For all sufficiently large \(M\), at the finite parameters
chosen above, there is a continuous rule \(\mathfrak b\) on
the uniform topology of the fields \((N,f,df,w)\) on a fixed
compact original subdomain, with values in a finite convex
hull in \(\mathcal B_M\), such that every limit solution at
\(t=1\) satisfies
\begin{equation}\label{eq:ah-contact-majorant}
 \mathfrak b(N,f,df,w)>N(H_g-zk)
\end{equation}
at every possible contact on \(S\) of a perimeter-minimizing
enclosure of \(S\).
\end{lemma}

\begin{proof}
Consider all limits at \(t=1\) arising from arbitrary smooth
choices \(b_S\in\mathcal B_M\) on arbitrary expanding
truncations.  Their restrictions to each fixed compact
region form a compact family in the uniform field topology.
Precompactness follows from Lemma~\ref{lem:ah-limit-mass};
closedness follows by a diagonal choice of actual truncation
solutions from the defining sequences of a convergent
sequence of limits.  This definition uses no existence
assertion, and the family may initially be empty.

The fixed-parameter end estimates are uniform over this
family.  Choose a large fixed radial sphere, beyond \(r_1\),
so that all spheres from slightly before it to infinity are
strictly outward mean convex in every limit metric.
Adjoin smooth compact orientable handlebodies to all components of
\(S\), possible because each is a closed orientable surface.  Choose
the gluing maps to reverse the induced boundary orientations, so the
orientation of \(\Omega\) extends over the caps.
Let \(O\) denote their union and \(X\) the resulting
one-ended smooth manifold without boundary; let \(K_c\)
be the compact region cut off by the chosen large sphere.
Extend each continuous positive metric through \(S\) to
the cap side, continuously with respect to its uniform
field values.  For example reflect in a fixed base collar
and interpolate to a fixed cap metric.  The mixed metric
entries vanish at \(S\), because \(f\) has constant trace.
Perimeters of sets containing \(O\) depend only on the
metric outside \(O\) and its boundary values.

For each of these metrics minimize perimeter among
finite-perimeter subsets of \(K_c\) containing \(O\)
up to null sets.  Compactness and lower semicontinuity
give a minimizer.  These facts hold for a continuous
positive metric by uniform smooth approximation of the
area integrands.  The sphere at \(r_1\) is an admissible
comparison, so every minimizer has perimeter at most
\(CR_s^2\).

Define contact points on \(S\) as points in the support
of a minimizer's perimeter measure.  The union of all
such contacts equals the contact set of one minimizer.
To prove this, minimize volume in an auxiliary fixed
smooth metric among the perimeter minimizers.  Perimeter
submodularity shows that the intersection and union of
any two minimizers are still minimizers.  The
volume-minimizing one, denoted \(E_-\), is therefore
contained almost everywhere in every other minimizer.
If a contact of a second minimizer were absent from the
perimeter support of \(E_-\), a small ball there would
have zero \(E_-\)-perimeter.  Since \(O\) occupies a
positive part of that ball, the ball would be filled
by \(E_-\), and hence by the second minimizer, a
contradiction.  The contact union is consequently the
closed contact set of \(E_-\).

These contact unions are upper semicontinuous in the
uniform field topology.  Indeed, along uniformly
convergent metrics, any sequence of minimizers has an
\(L^1\) subsequence converging to a minimizer for the
limit metric.  At a contact point, the complement has
uniform positive volume density: filling a chart ball
is an allowed competitor, so its perimeter in the ball
is bounded by a constant times its trace area on the
sphere.  Isoperimetry, applied to the complement inside
the ball, gives \(V^{2/3}\le CV'\), and thus
\(V(r)\ge cr^3\).  Constants are uniform on the compact
family of positive continuous metrics.  The cap supplies
positive density for the set on the other side.
These two density bounds persist when both the metric
and the contact point converge, so the limit point is
in the perimeter support of a limit minimizer.

The same density bounds place all contacts in the
measure-theoretic boundary.  Up to a set of
two-dimensional Hausdorff measure zero they therefore
belong to the reduced boundary, by the finite-perimeter
structure theorem.  At almost every point of their
intersection with the smooth surface \(S\), the tangents
are those of \(S\).  Since \(df\) is tangentially zero,
the area density there is \(x\,\dd A_g\).
Writing \(C_h\) for the contact union, we obtain
\begin{equation}\label{eq:ah-contact-area}
 \int_{C_h}x\,\dd A_g\le CR_s^2.
\end{equation}

Set \(\Theta=[N(H_g-zk)]_+\) on \(S\).
Since \(H_g\le k\), its positive part is bounded by \(Cs\).
For example, when \(k\ge0\),
\(N(H_g-zk)\le Nk(1-z)\le ks\); when \(k<0\)
the left side is nonpositive.  The upper lapse bound
also gives \(\Theta\le C\), with \(C\) independent of
the pointwise cap \(B_*\).  Fix \(B_*\) strictly above
this latter bound.  By \eqref{eq:ah-penalty-targets},
\(s\le x/T\) at \(S\), and hence
\begin{equation}\label{eq:ah-contact-flux-cost}
 \int_{C_h}\Theta^{p_1}\,\dd A_g
 \le C\int_{C_h}s\,\dd A_g
 \le C R_s^2/T=CM^{-1/5}.
\end{equation}
The admissible \(p_1\)-power budget in
\eqref{eq:ah-flux-class} is \(M^{-1/10}\).
Thus for sufficiently large \(M\) there is strict slack.
For each field in the compact family, approximate its
continuous nonnegative threshold on the closed set \(C_h\)
from above by a smooth function, with a small positive
margin, and taper it to zero in a sufficiently small
neighborhood of that set.  Regularity of surface measure
and \eqref{eq:ah-contact-flux-cost} keep its \(L^{p_1}\)
norm within the budget; the fixed gap below \(B_*\)
keeps its pointwise values admissible.  This constructs
a datum strictly majorizing the threshold on \(C_h\).
When \(C_h\) is empty the condition is vacuous.

Upper semicontinuity of contacts and continuity of the
threshold show that the same datum works in an open
neighborhood of the given field in the compact limit
family.  Take a finite such cover, with associated
smooth data \(b_1,\dots,b_m\), and enlarge the neighborhoods
to open subsets of the ambient uniform field space so
that their intersections with the limit family retain
this property.  A continuous partition of unity,
supplemented by the complement of the compact limit
family carrying datum zero, yields \(\mathfrak b\).
Explicitly, distance to the complements of these finitely
many open neighborhoods, together with distance to the
compact limit family, gives nonnegative continuous weights
with nonzero sum after a harmless subordinate refinement.
Normalize and take the corresponding convex combination.
At a field in the limit family the zero datum has weight
zero, and every active \(b_i\) is a strict majorant
at all its contacts.  Convexity preserves both constraints
in \eqref{eq:ah-flux-class}.  If the family is empty, use
the constant zero rule.  Only the original compact
subdomain's fields enter this rule; no trial fields on
the caps are required.
\end{proof}

\subsection{Existence at fixed finite parameters}

\begin{proposition}\label{prop:ah-finite-existence}
With the choices above, the system
\eqref{eq:ah-system}--\eqref{eq:ah-outer-data}, using
\(b_S=\mathfrak b(N,f,df,w)\), has a smooth positive
solution at \(t=1\) on every sufficiently large fixed
truncation.  Expanding truncations have a subsequence
converging to a solution smooth up to \(S\) on the
whole exterior.  It satisfies all the preceding estimates
and realizes the rule \(\mathfrak b\).
\end{proposition}

\begin{proof}
We describe the compact map precisely, because a frozen
source need not preserve positivity of a lapse output
away from fixed points.  Fix a small Hölder exponent
\(\gamma>0\), to be decreased below the exponents in
the a priori estimates, and use triples of
\(C^{1,\gamma}\) functions on the fixed truncation.
Choose \(a>0\) so that \(2a\) is below the fixed-point
floor in Lemma~\ref{lem:ah-lapse}, and a smooth function
\(\mathcal E:\R\to[a,\infty)\) equal to the identity
on \([2a,\infty)\).

Given trial fields and \(t\), first solve the linear
nondivergence height equation, freezing
\(\mathcal P,\beta,d\log\mathcal E(N)\) from the
trial arguments, with Dirichlet values \(tM\) and zero.
At zero trial gradient the smooth continuations are
\(\mathcal P=g^{-1}\) and \(\beta=2\).
The first-order coefficient acts on the new height.
Next compute \(v\) from this new height and the
extended trial lapse, and set \(Z=K(v)\).
Solve the linear density equation for the raw new lapse:
\[
 \divg_g(\nabla N_{\mathrm{new}}-ZN_{\mathrm{new}})
    =3N_{\mathrm{new}}+\Sigma_{\mathrm{trial}} .
\]
Use the prescribed outer current flux and inner current
flux \(-t\mathfrak b\) evaluated on the trial fields.
The source here is frozen from the trial fields,
including the raw trial lapse in its cutoff at
\(N\le\epsilon\).  The scalar multiplying the drift
and the zeroth term is the raw output lapse, not
\(\mathcal E(N_{\mathrm{new}})\).
The adjoint argument in Lemma~\ref{lem:ah-lapse}
proves unique linear solvability.

Finally solve a linear \(w\)-equation whose coefficient
\(A\) is formed from the new height and
\(\mathcal E(N_{\mathrm{new}})\).  Its right side is
\[
 3\mathcal E(N_{\mathrm{new}})
    +\Sigma_{\mathrm{trial}}x_{\mathrm{trial}}^+
    +P_{\mathrm{trial}},
 \qquad
 x_{\mathrm{trial}}^+
  =s(\mathcal E(N_{\mathrm{trial}}),df_{\mathrm{trial}})
      (w_{\mathrm{trial}})^+ .
\]
It has the zero inner natural flux and the prescribed
positive outer Dirichlet datum.  The three solves are
sequential, and define a continuous compact map
\(\mathcal T_t\).
Both frozen sources include their prescribed factor \(t\).
Indeed bounded sets of input fields give uniform
ellipticity and Hölder coefficients.  The height output
is bounded in \(C^{2,\gamma}\), so the drift in the
lapse step is \(C^{1,\gamma}\); its inverse estimates
are uniform on bounded input sequences by Schauder
theory, coefficient compactness in a lower exponent,
and the zero-kernel adjoint argument.  The lapse output
and then the \(w\) output have second-order Hölder
bounds.  The rule \(\mathfrak b\) has values in a
fixed finite convex hull of smooth functions, hence
uniform bounds of every spatial derivative.  The
compact embedding of these \(C^{2,\gamma}\) bounds
into \(C^{1,\gamma}\) proves compactness.  The same estimates and
uniqueness prove continuity without differentiating
the selection rule.

At a fixed point, the cutoff in the actual lapse
variable and the bounded drift give the positive floor
in Lemma~\ref{lem:ah-lapse}.  Thus every extension
\(\mathcal E\) becomes the identity.  The last equation
and positive outer datum give \(w>0\), so the positive
part also disappears.  Every fixed point therefore
satisfies the original equations and boundary data.

All such fixed points for \(0\le t\le1\) have bounded
\(C^{1,\gamma}\) norms for the fixed truncation and
fixed parameters.  Here is the regularity chain
underlying this assertion.  The lapse bound, height
gradient bound, and bounded sources give lapse
\(W^{1,q}\) estimates, \(q>3\), and Hölder estimates
for the height gradient, as in
Lemma~\ref{lem:ah-graph-estimates}.  The lapse divergence
equation then gives \(C^{1,\gamma'}\) estimates for
some \(\gamma'>0\); its scalar-forcing portion can
equivalently be estimated by the Laplace \(W^{2,q}\)
theorem.  The selected boundary data have bounded
smooth norms.  The nondivergence height equation
gives \(C^{2,\gamma'}\) estimates, after decreasing
\(\gamma'\) if necessary.  The \(w\)-equation has
bounded forcing, since \(\Sigma x\le L\) and \(P\)
has finite strength.  Mixed-boundary coercivity first
gives an \(H^1\) bound, then scalar local boundedness
and Hölder estimates give a uniform bound and
continuity.  Successive Schauder estimates now give
second-order Hölder bounds for all fields.  Take
the Banach exponent \(\gamma<\gamma'\).
The same bootstrap makes every fixed point smooth.

At \(t=0\), the map is constant.  The first output
is \(f_0=0\), so the drift in the second step is
zero independently of the input.  Its output is
the unique positive solution \(N_0\) of
\(\Delta_gN_0=3N_0\), with zero inner flux and the
prescribed outer flux.  The third output \(w_0\)
uses \(\mathcal E(N_0)=N_0\), by the homogeneous lapse
floor, and is the unique solution of
\[
 -\divg_g(N_0^2\nabla w_0)=3N_0,\qquad
 \partial_\nu w_0|_S=0,\qquad w_0|_{S_R}=1/V_0 .
\]
In general \(w_0\ne1/N_0\), because the outer lapse
condition is Neumann.  Choose a ball in the Banach
space containing this constant output and all
homotopy fixed points.  On its boundary
\(\Id-\mathcal T_t\) never vanishes, and its
Leray--Schauder degree is \(+1\) at \(t=0\).
Homotopy invariance gives a fixed point at \(t=1\).

Now let \(R\to\infty\).  The fixed-parameter compact
estimates give a subsequential limit on every compact
set, and a diagonal subsequence gives a global limit.
It belongs to the family used to define \(\mathfrak b\).
Uniform convergence of \((N,f,df,w)\) on its defining
compact domain implies convergence of the selected
datum, in every spatial norm because its range is
a finite smooth convex hull.  Thus the limit realizes
the rule, and the elliptic bootstrap is valid up to
\(S\) as well as in the interior.
\end{proof}

\subsection{A free minimal boundary}

Apply Proposition~\ref{prop:ah-finite-existence} at \(t=1\).
For this smooth metric, extend through the cap side
smoothly, and consider the perimeter minimization in
Lemma~\ref{lem:ah-contact-selector}.  Equation
\eqref{eq:ah-boundary-curvature} and the strict majorant
give \(H_h(S)<0\) at every possible contact.
This excludes contact without any assumed regularity
of the minimizing surface at the obstacle.

Here is the finite-perimeter argument.  Near the closed
contact set, the unit normals to sufficiently small
parallel outward collars of \(O\) have negative
divergence.  Off a neighborhood of that set, a thin
collar is already filled up to null sets: there is no
perimeter support on \(S\), and the cap side is filled.
Adjoin the entire collar of width \(a\) to a minimizer.
For almost every sufficiently small \(a\), Gauss--Green
and slicing bound the perimeter change above by the
integral of that negative divergence on the added part.
The comparison field has unit flux on the new outer
leaf, and its flux across the removed old boundary is
bounded by the area removed.  If contact existed,
the complement density proved above would give
positive added volume, so the perimeter would strictly
decrease.  This is impossible.

There is likewise no contact with the outer constraint
\(\partial K_c\).  The outward normals to the far
coordinate spheres have positive divergence, so
cutting away the outside part of a competitor decreases
perimeter whenever it removes positive volume.
Gauss--Green and slicing make this statement valid
for finite-perimeter competitors at almost every
cutting radius.  It also shows that no larger precompact
competitor, unconstrained by \(K_c\), can have smaller
perimeter: first trim it by such a sphere.

All minimizing boundaries are therefore free, locally
perimeter minimizing in a smooth three-dimensional
metric.  Interior codimension-one minimizing-boundary
regularity gives smooth compact embedded minimal
boundaries.  Choose a largest minimizer \(E\), by
maximizing auxiliary volume among the minimizers in
\(K_c\).  Its open representative contains \(O\)
compactly in its interior and is strictly outward
minimizing: equality for an enclosure differing by
positive volume would
contradict volume maximality after trimming.
Every component of \(E\) meets \(O\), since an
unattached component could be removed with a strict
perimeter saving.  Its open exterior is connected,
since filling any bounded exterior component would
also save perimeter.

Write
\begin{equation}\label{eq:ah-hull-area}
 A_h=\Area_h(\partial E)>0.
\end{equation}
The closed exterior \(X\setminus E\), with its boundary
included, is smooth, connected, complete, orientable, and one-ended.
Its orientation is the restriction of the orientation of \(X\);
the graph and conformal changes alter the metric, not this orientation.
Its compact truncations are compact, and its end is
asymptotically hyperbolic.  It satisfies
\(\Scal_h\ge-6\) by \eqref{eq:ah-scalar}, and its
weighted scalar integrability and finite mass fluxes
were verified in Lemma~\ref{lem:ah-limit-mass}.
Its minimum enclosing area is exactly \(A_h\).
Indeed any admissible end-side domain in this closed
exterior, on taking the complement of its interior
in \(X\), supplies a precompact finite-perimeter
enclosure of \(E\).  Its perimeter is bounded by the
entire intrinsic boundary area specified in the
definition, including any part coinciding with
\(\partial E\).  Outward minimization gives the
lower bound \(A_h\), and the original boundary
itself gives equality.

The time-symmetric hypothesis of
Theorem~\ref{thm:ah-reduction} is thus applicable to
this fixed metric and gives
\begin{equation}\label{eq:ah-ts-application}
 p_0(h)\ge F(A_h)\ge0.
\end{equation}
No causal property of the intermediate mass covector
has been used.

\subsection{The weighted bad-volume estimate and conclusion}

Suppose first that \(A_*>0\), fix \(0<\varepsilon_*<1\),
and choose the preceding \(\delta\) sufficiently small
for Proposition~\ref{prop:area-transfer}.
Fix a sufficiently large projection radius \(r_p\)
and put
\[
 \rho=\min(1,r_p^2/r^2),\qquad
 \mathcal D=\{x<1-\delta\}.
\]
By \eqref{eq:ah-mass-budget} and
\eqref{eq:ah-ts-application},
\[
 \int_\Omega(1-x)\Sigma
 =\int_S b_S+\int_\Omega P
       -8\pi\{p_0(h)-p_0(g)\}\le C.
\]
The constant is independent of \(M\).
Lemma~\ref{lem:ah-penalties} makes
\(\Sigma=L\) on \(\mathcal D\cap\{r\le R_s\}\),
and \(1-x\ge\delta\) there.  Consequently
\[
 \Vol_g(\mathcal D\cap\{r\le R_s\})\le C/L.
\]
Since \(W\le1+\sqrt{W^2-1}\), Cauchy--Schwarz and
\eqref{eq:ah-height-energy} give an inner weighted
volume bound \(C(L^{-1}+\sqrt{M/L})\).
For large \(M\), \(R_s>r_p\).  The asymptotically
hyperbolic volume element gives
\[
 \int_{\{r>R_s\}}\rho^{3/2}\,\dd V_g\le C/R_s,\qquad
 \int_{\{r>R_s\}}\rho^3\,\dd V_g\le C/R_s^4 ,
\]
where \(r_p\) is fixed and absorbed into \(C\).
A second application of Cauchy--Schwarz yields
\begin{equation}\label{eq:ah-bad-volume}
 \int_{\mathcal D}\rho^{3/2}W\,\dd V_g
 \le C\left(L^{-1}+\sqrt{M/L}
                  +R_s^{-1}+\frac{\sqrt M}{R_s^2}\right)
 \longrightarrow0 .
\end{equation}
For the choices \eqref{eq:ah-parameter-scales},
all four powers tend to zero.

We check the hypotheses of the area-transfer proposition
explicitly.  The capped smooth manifold \(X\) and cap
set \(O\) are fixed, so the relevant second Betti number
of \(X\setminus\operatorname{int}O\) is finite and fixed.
The chosen \(E\) is a cap-only largest minimizing
enclosure, has a free smooth compact minimal boundary,
contains \(O\) compactly, is strictly outward minimizing,
and each component meets \(O\).  The comparison metric
is the smooth \(h=x\bar g\); here there is no seam or
smoothing error.  It obeys
\(\bar g\ge g\), \(\dd V_{\bar g}=W\,\dd V_g\), and
\(\Scal_h\ge-6\) outside \(E\).
Its asymptotically hyperbolic end supplies the proper
subsolution required in Proposition~\ref{prop:area-transfer}.
Lemma~\ref{lem:area-projection} gives, on every relevant
enclosing level with connected exterior,
\[
 \int\rho\,\dd A_g\ge(1-o_{r_p}(1))A_*.
\]
Choose \(r_p\) first to make this error sufficiently
small for the fixed \(\varepsilon_*\), then choose
\(\delta\), and finally take \(M\) large.
The proposition's temporary conformal modification
for starting the flow may be arbitrarily small at
each fixed metric and leaves the initial area
unchanged.  Its uniform scalar lower bound therefore
has the required fixed constant, for example \(7\).
Equation~\eqref{eq:ah-bad-volume} meets the remaining
smallness hypothesis.  Proposition~\ref{prop:area-transfer}
gives
\begin{equation}\label{eq:ah-area-conclusion}
 A_h\ge(1-\varepsilon_*)A_* .
\end{equation}

Combine \eqref{eq:ah-mass-budget},
\eqref{eq:ah-ts-application}, and
\eqref{eq:ah-area-conclusion}.  For every fixed
\(\varepsilon_*>0\) and arbitrarily large \(M\),
\[
 p_0(g)+o(1)\ge F\bigl((1-\varepsilon_*)A_*\bigr).
\]
First let \(M\to\infty\), then
\(\varepsilon_*\downarrow0\), using continuity of \(F\).
If \(A_*=0\), no area transfer is needed:
\(F(A_h)\ge F(0)\) and the mass budget give the
same conclusion directly.  This proves
Theorem~\ref{thm:ah-reduction}.
